% =====================================================================
%  Sezione 2.9 — Selezione delle grandezze tramite DTW
%  Sostituisce la sezione "Analisi di separabilita delle feature con DTW".
% =====================================================================

\section{Selezione delle grandezze tramite \emph{Dynamic Time Warping}}
\label{sec:dtw}

Costruito il bacino di trentadue grandezze candidate, resta da stabilire quali
diano effettivamente una firma temporale distinguibile fra caduta e attività
quotidiana. Una grandezza può essere geometricamente ben definita e
computazionalmente economica senza per questo separare le due classi.

\subsection{Perché il DTW}

Il confronto avviene fra serie temporali di durata e velocità diverse: due
cadute possono avere la stessa struttura dinamica e svolgersi in tempi
differenti, oppure presentare il picco di una grandezza qualche fotogramma prima
o dopo. Una distanza punto per punto, come quella euclidea, giudicherebbe
diverse due curve che hanno la stessa forma ma sono sfasate.

Il \emph{Dynamic Time Warping} confronta invece la \emph{forma} delle due serie,
deformando l'asse dei tempi alla ricerca dell'allineamento di costo minimo. È
per questo lo strumento adatto a valutare se una grandezza abbia un andamento
temporale coerente all'interno di una classe e diverso fra classi.

\subsection{Le tre scelte metodologiche}

\paragraph{Normalizzazione per curva.} Le grandezze hanno scale eterogenee —
gradi, velocità, rapporti adimensionali — e sui valori grezzi la distanza
sarebbe dominata da quelle di ampiezza maggiore. Ogni curva viene quindi
riportata a media nulla e deviazione unitaria. Si perde l'informazione di
ampiezza assoluta, il che è accettabile perché in questa fase interessa la
forma.

\paragraph{Vincolo sul cammino di allineamento.} Un DTW libero può deformare
l'asse dei tempi al punto da far somigliare la curva di un'attività normale a
quella di una caduta. Si impone quindi che l'allineamento fra due indici
temporali resti vicino alla diagonale, $|i-j| \leq w$ con $w = 10$ fotogrammi:
sono ammessi piccoli sfasamenti, non deformazioni prive di senso fisico.

\paragraph{Ricampionamento a lunghezza comune.} È una scelta che deriva
direttamente dalla precedente, e la sua necessità è stata individuata in fase di
verifica. Il vincolo $|i-j| \leq w$ non è applicabile a due curve di lunghezza
diversa senza allargarlo ad almeno la loro differenza, altrimenti nessun cammino
di allineamento raggiunge il termine. Su questo dataset le riprese di caduta
durano circa 180 fotogrammi e quelle di attività quotidiana circa 820: le coppie
\emph{interne} a ciascuna classe hanno quindi lunghezze simili e mantengono la
banda stretta, mentre le coppie \emph{incrociate} — proprio quelle che devono
risultare distanti — se la vedono allargare a oltre seicento e si allineano
quasi senza vincolo.

L'effetto è sistematico e non uniforme: la distanza fra le classi risulta
deflazionata, e la deflazione colpisce selettivamente le grandezze integrali,
che sono curve regolari e sotto allineamento libero si deformano l'una
nell'altra a costo quasi nullo. Prima della correzione, cinque delle nove
grandezze integrali risultavano al di sotto della soglia di separazione e
venivano scartate come rumore.

Portando ogni curva a una lunghezza comune il problema si annulla, perché tutte
le coppie hanno la stessa lunghezza e la banda vale identica per tutte. Si perde
l'informazione di durata, ed è accettabile per una ragione precisa: la durata di
una ripresa è identica per tutte e trentadue le grandezze, quindi non può
contribuire a distinguerle fra loro, mentre il rumore che introduce sì.

\subsection{Il punteggio di separabilità}

Per ciascuna grandezza si calcolano tre popolazioni di distanze: fra curve di
caduta ($D_{FF}$), fra curve di attività quotidiana ($D_{AA}$) e incrociate
($D_{FA}$). Una grandezza discriminativa produce curve simili all'interno della
stessa classe e diverse fra classi, quindi $D_{FF}$ e $D_{AA}$ bassi e $D_{FA}$
alto. Il punteggio è

\[
S = \frac{\mu_{FA}}{\mu_{\text{intra}} + \varepsilon},
\qquad
\mu_{\text{intra}} = \frac{\mu_{FF} + \mu_{AA}}{2},
\]

affiancato da una versione che tiene conto anche della dispersione,

\[
S_{\text{rob}} = \frac{\mu_{FA} - \mu_{\text{intra}}}
                      {\sigma_{FA} + \sigma_{\text{intra}} + \varepsilon},
\]

perché un punteggio alto può derivare da valori anomali o da forte variabilità
invece che da una separazione reale. Un valore $S = 1$ indica assenza di
separazione: le curve di classi diverse distano quanto quelle della stessa
classe.

\subsection{Il criterio del minimo fra le due viste}

L'analisi iniziale calcolava un punteggio unico, mescolando le riprese delle due
telecamere. Questa scelta nasconde però il divario più grande misurato in tutto
il lavoro.

UP-Fall riprende ogni evento con due telecamere simultanee: una laterale,
all'altezza del busto, e una frontale rialzata. Una telecamera che guarda
dall'alto comprime la rotazione del tronco lungo il proprio asse ottico: la
stessa grandezza che di profilo mostra ottantaquattro gradi di escursione fra
persona eretta e distesa ne mostra tre dall'alto. Una grandezza che separa
benissimo da una vista può quindi non separare nulla dall'altra, e mescolando le
due il crollo resta invisibile fino alla misura sul modello finito.

Il punteggio viene perciò calcolato \emph{separatamente per telecamera}, e la
selezione usa il \textbf{minimo} fra i due:

\[
S_{\text{robusto}} = \min\big(S_{\text{cam}_1},\ S_{\text{cam}_2}\big).
\]

Il minimo e non la media, perché una grandezza ottima da una vista e inutile
dall'altra è esattamente quella da escludere, e la media la premierebbe. Il
costo di questo criterio va dichiarato: scarta grandezze valide da una vista
sola. L'asse maggiore dell'ellisse, per esempio, ottiene sulla vista dall'alto
il secondo punteggio più alto dell'intera analisi ma un valore mediocre di
profilo, e viene escluso. È il prezzo di un sistema che deve funzionare senza
sapere come sia stata montata la telecamera.

\subsection{Ridondanza fra le grandezze scelte}

Il DTW valuta ogni grandezza \emph{isolatamente}: non sa nulla delle altre, e
quindi se lo stesso fenomeno fisico è misurato in tre modi lo premia tre volte.

Sui dati questa non è un'eventualità teorica. Le prime tre grandezze della
classifica risultavano correlate fra loro al di sopra di $0{,}94$, ed erano
tutte descrizioni di «quanto è verticale il tronco»: l'altezza della testa
rispetto al bacino in due rappresentazioni diverse e l'angolo del tronco. Il
legame è geometrico — il tronco è un segmento rigido, quindi l'altezza della
testa vale approssimativamente $L\cos\theta$ — e prenderle tutte e tre
significherebbe occupare tre posti con un segnale solo.

La selezione è quindi resa consapevole della ridondanza: si scorre la classifica
dall'alto e si aggiunge una grandezza soltanto se la sua correlazione assoluta
con \emph{tutte} le già scelte resta al di sotto di una soglia, fissata a
$0{,}85$. Le grandezze escluse vengono dichiarate insieme a quella che le
blocca, così l'esclusione è documentata invece di essere semplicemente
un'assenza dall'elenco.

Un'ultima cautela riguarda i fotogrammi privi di posa: la correlazione si misura
sui soli fotogrammi con posa valida, perché sugli altri tutte le grandezze
valgono zero per costruzione e risulterebbero correlate fra loro per un motivo
che non ha nulla a che vedere con la dinamica del corpo.

\subsection{Un limite dichiarato}

L'analisi qui descritta è una \textbf{preselezione euristica}, non una
dimostrazione. Nessuna grandezza, presa singolarmente, raggiunge un punteggio
tale da separare le classi da sola: i valori si collocano fra $1{,}05$ e
$1{,}31$, mentre $1{,}0$ indica assenza di separazione. Il sottoinsieme va
valutato \emph{insieme}, e l'arbitro è la validazione incrociata sul modello,
riportata nel Capitolo~\ref{cap:sperimentale}.

Va inoltre osservato che la selezione dipende dal campione: su un pilota a due
soggetti la classifica risultava sensibilmente diversa da quella ottenuta su
dieci, e due grandezze che apparivano fra le migliori si rivelavano al di sotto
della soglia di separazione una volta esteso il campione.
